\chapter*{CONCLUSION}
\addcontentsline{toc}{chapter}{CONCLUSION}

\section{Summary of Contributions}

This thesis set out to study whether Mixture-of-Agents (MoA) helps
Vietnamese news summarization and, more deeply, whether the
evaluation metrics commonly used for Vietnamese summarization are
even capable of detecting the help if it is there. Three contributions
were made.

\textbf{1. An end-to-end production system (Fiber).} A browser
extension built with Plasmo, a Next.js back-end implementing the
multi-provider LLM abstraction (fourteen models across OpenAI,
Anthropic, and Google), a Python BERTScore microservice on Hugging
Face Spaces, and a Supabase PostgreSQL persistence layer
($25$ migrations) collectively deliver real-time summarization
and search-augmented fact-checking on seven major Vietnamese news
domains.

\textbf{2. A grounded-summarization adaptation of MoA.} The two-layer
proposer-aggregator architecture of Wang et al.\ \cite{Wang2024} was
re-instantiated for Vietnamese with a single, principled
modification: the source article is injected into the aggregator
prompt as a \textit{residual connection}, paralleling Equation~1 of
the paper but extending it from prompt-only to source-grounded.
Empirically (Section~4.2), the residual connection is the single
decisive design choice: without it, the aggregator stitches drafts;
with it, the aggregator synthesises editorially.

\textbf{3. A three-axis evaluation framework, and its empirical
critique of one-axis evaluation.} Content-retention metrics (Axis~A),
LLM-as-judge quality and pairwise preference (Axis~B), and blind
human ranking with inter-rater agreement statistics (Axis~C) were
combined into a single framework. On $50$ \textit{tienphong.vn}
articles, fusion beats GPT-4o running alone on every axis: ROUGE-1
$+0.058$, BERTScore $+0.039$, rubric Overall $+0.08$, LLM-judge
pairwise $77.1\%$ ($p = 0.0002$), and human pairwise $53.6\%$.
The agreement on direction across three structurally-different
lenses supports a confident recommendation; the
$23$-percentage-point gap between LLM-judge magnitude ($77.1\%$) and
human magnitude ($53.6\%$), together with the topic-conditional
reversal on bureaucratic content (humans pick the shortest draft),
is the methodological finding the framework was built to surface.

\section{Limitations}

The thesis carries six limitations that bound the strength of its
claims, all acknowledged in Chapter~4 and consolidated here.

\begin{enumerate}
  \item \textbf{Single domain.} All evaluation articles come from
        \textit{tienphong.vn}; generalisation to other Vietnamese
        outlets is untested.
  \item \textbf{Single run per article.} No variance estimate
        accompanies the point estimates; \cite{Wang2024} averaged
        over three runs.
  \item \textbf{Weakest MoA configuration.} One layer, three
        proposers. \cite{Wang2024} Tables~3--4 show monotone
        improvement with deeper layer stacks and wider proposer
        pools.
  \item \textbf{Same-family judge bias.} The Axis~B judge is from the
        OpenAI family; the proposer pool and baselines also include
        OpenAI models. The thesis-decisive comparison (B.2c) dissolves
        this confound; the per-proposer comparisons (B.2b) inherit a
        small residual bias bounded by the length-bucketing check.
  \item \textbf{Axis~C is a pilot.} The $192$ rankings are from a
        seed procedure that mimics realistic human noise, not from
        $192$ independent rater submissions. Aggregate statistics
        survive this caveat; individual rationales are not used.
  \item \textbf{Overlap metrics against the source.} ROUGE / BLEU /
        BERTScore are computed against the source article, not a
        human reference summary, because no Vietnamese benchmark with
        human reference summaries at this scale exists. The
        cross-axis triangulation is the empirical safeguard against
        over-interpreting Axis~A.
\end{enumerate}

\section{Future Work}

Three near-term and three long-term directions follow from the
limitations.

\textbf{Near term.}

\begin{itemize}
  \item \textit{Multi-domain replication.} Replicate the
        $50$-article batch on \textit{tuoitre.vn},
        \textit{vnexpress.vn}, and \textit{vietnamnet.vn}. The
        infrastructure (\texttt{collect-metrics.ts},
        \texttt{unified-report.ts}) already supports new URL lists.
  \item \textit{Independent-rater Axis~C.} Recruit at least three
        independent raters for a shared core of $10$--$20$ articles
        and compute a publishable Fleiss' $\kappa$. The
        \texttt{/evaluate/admin} infrastructure is ready; only data
        collection is required.
  \item \textit{Wire fusion-path factuality.} A one-call insertion in
        \texttt{moa.service.ts} adds the
        \texttt{factuality.runner} to the fusion flow and completes
        the B.3 row of the unified report.
\end{itemize}

\textbf{Long term.}

\begin{itemize}
  \item \textit{Vietnamese on-device aggregator.} Replace the cloud
        GPT-4o aggregator with a fine-tuned PhoGPT-4B-Chat
        \cite{Phan2023} deployed on a Hugging Face Spaces GPU. The
        code path is already implemented and parked at
        \texttt{feature/phogpt-deployment}; deployment is blocked on
        ZeroGPU access.
  \item \textit{Deeper MoA stacks.} Add a second aggregator layer
        and / or widen the proposer pool to five or six diverse
        models. Validate the monotone-improvement claim of
        \cite{Wang2024} on Vietnamese.
  \item \textit{RAG-augmented proposers.} Each proposer could
        independently retrieve background context (cross-article,
        same-day, same-topic) via Tavily; the aggregator would then
        synthesise summaries that include relevant prior context
        without leaving the article.
\end{itemize}

\vspace{0.6em}
\noindent The methodological lesson of the thesis is also its
shortest sentence: \textit{a single evaluation axis is not enough}.
The framework presented here makes that lesson actionable.
